% s-240-p240-m002; original PDF 240, 241
\Aphorism{s-240-p240-m002}{333}
So wie die getheilte Strecke länger erscheint als die gleichlange ungetheilte, der geknotete Faden länger als der gleichlange glatte, so rechnen wir auch im geistigen Erfassen unseren inneren Receptionswiderstand dem Object zu: ein Buch in fremder Sprache dünkt uns geistreicher und bedeutender als ein deutsches Buch gleichen Niveaus, den schwerverständlichen Autor nehmen wir für vollwichtiger als den klaren und durchsichtigen. Unsere Art zu lesen ist geradezu — moralisch: die Ansprüche und Anmassungen, die sich der Verfasser gegen uns erlaubt, stellen wir noch in sein Guthaben, statt ins „Soll“! Auf diese unsere Gutmüthigkeit und verkehrte Buchung haben die meisten deutschen Philosophen gerechnet; Kant nicht ausgenommen, der keine kleine Angst hatte, zu leicht verstanden zu werden. Diese Philosophen haben Recht; man muss wirklich ein dummer Teufel von Autor-Ehrlichkeit sein, um dem Publicum, das benebelt und übervortheilt sein \Sperr{will}, reinen Wein einzuschänken. Ihr verlangt Mysterien; also — mystificiren wir euch!
\par

